\begin{explanationbox}
	\textbf{\textcolor{CExplanation}{一句话直觉}}

	集合运算就是把“元素是否属于集合”翻译成“或、且、非”三种逻辑条件，再据此筛选元素。

	\medskip
	\textbf{\textcolor{CExplanation}{核心拆解}}
	\begin{description}[style=nextline, leftmargin=2.8em, labelsep=0.5em, itemsep=0.45em]
		\item[\textbf{要点一（并集对应“或”）：}]
		      \(x\in A\cup B\) 表示 \(x\in A\) 或 \(x\in B\)。这里的“或”是相容或：只要至少一个条件成立即可，两个条件同时成立也包含在内。

		\item[\textbf{要点二（交集对应“且”）：}]
		      \(x\in A\cap B\) 表示 \(x\in A\) 且 \(x\in B\)，两个条件必须同时成立。

		\item[\textbf{要点三（补集对应“全集内取反”）：}]
		      \(x\in\complement_U A\) 表示 \(x\in U\) 且 \(x\notin A\)。补集依赖全集；若上下文已经固定全集，可以省略下标，但全集本身不能缺失。

		\item[\textbf{要点四（差集具有方向性）：}]
		      \(x\in A\setminus B\) 表示 \(x\in A\) 且 \(x\notin B\)。一般来说，\(A\setminus B\neq B\setminus A\)。
	\end{description}

	\medskip
	\textbf{\textcolor{CExplanation}{韦恩图直觉}}

	\(A\cup B\) 是两个区域的全部；\(A\cap B\) 是两个区域的重叠部分；\(\complement_U A\) 是全集内、集合 \(A\) 外的部分；\(A\setminus B\) 是 \(A\) 中删去与 \(B\) 重叠后剩下的部分。

	\medskip
	\textbf{\textcolor{CExplanation}{常用等价关系}}
	\begin{itemize}[leftmargin=*, itemsep=0.5em, topsep=0.4em]
		\item \(\complement_U A=U\setminus A\)：补集就是全集减去 \(A\)。
		\item \(A\setminus B=A\cap(\complement_U B)\)：在 \(A\) 中排除 \(B\)，等价于“属于 \(A\) 且属于 \(B\) 的补集”。
		\item \(A\setminus B=\varnothing\iff A\subseteq B\)：\(A\) 中没有落在 \(B\) 外面的元素，等价于 \(A\) 的所有元素都属于 \(B\)。
	\end{itemize}

	\medskip
	\textbf{\textcolor{CExplanation}{考点价值}}
	\begin{itemize}[leftmargin=*, itemsep=0.5em, topsep=0.4em]
		\item \textbf{直接运算：}给定列举集合、区间或数轴，求并集、交集、补集和差集。
		\item \textbf{参数问题：}把集合运算结果转化为包含关系、端点关系或不等式，再求参数范围。
		\item \textbf{条件转化：}把复合条件写成集合语言，例如“至少满足一个条件”对应并集，“同时满足”对应交集。
	\end{itemize}

	\medskip
	\textbf{\textcolor{CExplanation}{顿悟点}}

	\textbf{不要先背区域，要先判断任意元素 \(x\) 应满足什么条件；元素条件写对，集合运算就不会错。}
\end{explanationbox}